% Diagrama TikZ: Fluxo Integrado dos Processos Metodológicos
% Requer \usepackage{tikz} e \usetikzlibrary{...} no preâmbulo

\begin{tikzpicture}[
  % ── Estilos de nós ──────────────────────────────────────
  phase/.style={
    draw=#1!70!black, fill=#1!12, rounded corners=3pt,
    font=\scriptsize, text width=2.4cm, align=center,
    minimum height=2.0cm, inner sep=3pt, line width=0.8pt},
  phasenum/.style={
    circle, draw=#1!70!black, fill=#1!25,
    font=\scriptsize\bfseries, inner sep=1.5pt,
    minimum size=0.5cm},
  output/.style={
    draw=gray!60, fill=gray!6, rounded corners=2pt,
    font=\tiny\itshape, text width=2.4cm, align=center,
    minimum height=0.9cm, inner sep=2pt, dashed},
  method/.style={
    draw=gray!40, fill=white, rounded corners=1pt,
    font=\tiny, text width=2.4cm, align=center,
    minimum height=0.7cm, inner sep=2pt},
  titlebox/.style={
    draw=black!70, fill=black!7, rounded corners=3pt,
    font=\footnotesize\bfseries, text width=13.8cm, align=center,
    minimum height=0.8cm, inner sep=4pt},
  integbox/.style={
    draw=black!70, fill=black!7, rounded corners=3pt,
    font=\scriptsize, text width=13.8cm, align=center,
    minimum height=0.9cm, inner sep=4pt},
  arr/.style={-{Stealth[length=5pt]}, thick, draw=gray!60},
  dataarr/.style={-{Stealth[length=4pt]}, semithick, draw=blue!40, densely dashed},
]

% ══════════════════════════════════════════════════════════
% TÍTULO
% ══════════════════════════════════════════════════════════
\node[titlebox] (TITLE) {\textsc{Arquitetura Metodológica Integrada}};

% ══════════════════════════════════════════════════════════
% FASES (linha horizontal)
% ══════════════════════════════════════════════════════════
\node[phase=blue, below=1.0cm of TITLE.south west, anchor=north west] (F1) {
  \textbf{Fase I}\\[2pt]
  Revisão de Escopo\\PRISMA-ScR\\[2pt]
  {\tiny Meta-análise de acurácia\\Conformidade FAIR\\Validação espacial}};

\node[phase=teal, right=0.3cm of F1] (F2) {
  \textbf{Fase II}\\[2pt]
  Adaptação Transcultural\\+ Delphi\\[2pt]
  {\tiny WOCAT-SLM-QBR\\Elicitação estruturada\\Variáveis linguísticas}};

\node[phase=violet, right=0.3cm of F2] (F3) {
  \textbf{Fase III}\\[2pt]
  Modelagem Fuzzy\\TRI + Fuzzy $\rightarrow$ ISB\\[2pt]
  {\tiny Funções de pertinência\\Regras de inferência\\Defuzzificação}};

\node[phase=orange, right=0.3cm of F3] (F4) {
  \textbf{Fase IV}\\[2pt]
  Auditoria Computacional\\Aceitação Institucional\\[2pt]
  {\tiny Completude documental\\Painel regulatório\\Conformidade FAIR}};

\node[phase=green!60!black, right=0.3cm of F4] (F5) {
  \textbf{Fase V}\\[2pt]
  Design Participativo\\Validação Territorial\\[2pt]
  {\tiny Co-design comunitário\\Usabilidade (SUS)\\Adoção continuada}};

% ══════════════════════════════════════════════════════════
% OUTPUTS
% ══════════════════════════════════════════════════════════
\node[output, below=0.5cm of F1] (O1) {Diagnóstico de prontidão\\operacional do ML\\para SAT};
\node[output, below=0.5cm of F2] (O2) {Instrumento\\WOCAT-SLM-QBR\\+ variáveis validadas};
\node[output, below=0.5cm of F3] (O3) {Índice de Salvaguarda\\Biocultural\\(ISB: 0--100)};
\node[output, below=0.5cm of F4] (O4) {Taxa de aceitação\\institucional e\\conformidade regulatória};
\node[output, below=0.5cm of F5] (O5) {Sistema validado\\territorialmente\\com indicadores SUS};

% ══════════════════════════════════════════════════════════
% SETAS ENTRE FASES (sequencial)
% ══════════════════════════════════════════════════════════
\draw[arr] (F1.east) -- (F2.west);
\draw[arr] (F2.east) -- (F3.west);
\draw[arr] (F3.east) -- (F4.west);
\draw[arr] (F4.east) -- (F5.west);

% Setas verticais fase → output
\foreach \i in {1,...,5}{
  \draw[arr] (F\i.south) -- (O\i.north);
}

% ══════════════════════════════════════════════════════════
% SETAS DE RETROALIMENTAÇÃO (dados entre fases)
% ══════════════════════════════════════════════════════════
\draw[dataarr] (O1.south) -- ++(0,-0.3) -| 
  node[pos=0.25, below, font=\tiny\itshape, text=blue!50]{lacunas $\rightarrow$ variáveis} (F2.south);
\draw[dataarr] (O2.south) -- ++(0,-0.6) -|
  node[pos=0.25, below, font=\tiny\itshape, text=blue!50]{variáveis $\rightarrow$ fuzzy} (F3.south);
\draw[dataarr] (O3.south) -- ++(0,-0.3) -|
  node[pos=0.25, below, font=\tiny\itshape, text=blue!50]{ISB $\rightarrow$ dossi\^e} (F4.south);

% ══════════════════════════════════════════════════════════
% CAPÍTULOS (mapeamento)
% ══════════════════════════════════════════════════════════
\node[font=\tiny\bfseries, text=gray!50, below=0.05cm of O1] {Cap.\,1};
\node[font=\tiny\bfseries, text=gray!50, below=0.05cm of O2] {Cap.\,2};
\node[font=\tiny\bfseries, text=gray!50, below=0.05cm of O3] {Cap.\,3};
\node[font=\tiny\bfseries, text=gray!50, below=0.05cm of O4] {Cap.\,4};
\node[font=\tiny\bfseries, text=gray!50, below=0.05cm of O5] {Cap.\,5};

% ══════════════════════════════════════════════════════════
% INTEGRAÇÃO FINAL
% ══════════════════════════════════════════════════════════
\node[integbox, below=1.6cm of O3] (INT) {
  \textbf{\textsc{Integração}} ---
  Sensoriamento remoto · Governança de dados (FAIR) · Inferência fuzzy ·
  Métricas auditáveis de salvaguarda · Proteção de PI};

\foreach \i in {1,...,5}{
  \draw[arr, gray!40] (O\i.south) -- ++(0,-0.5) -| (INT.north);
}

\end{tikzpicture}
